% ============================================================
%  Probeklausur Template (my_dhbw Stil, klausur_*.tex)
%  Formell mit Deckblatt-Header; Lösungen als grüne tcolorbox.
%  Renderer: pdflatex (zweimal)
% ============================================================
\documentclass[11pt, a4paper]{article}

\usepackage[ngerman]{babel}
\usepackage{fontspec}
\setmainfont[
  Path=/usr/share/fonts/TTF/,
  Extension=.ttf,
  BoldFont=DejaVuSans-Bold,
  ItalicFont=DejaVuSans-Oblique,
  BoldItalicFont=DejaVuSans-BoldOblique
]{DejaVuSans}
\setsansfont[
  Path=/usr/share/fonts/TTF/,
  Extension=.ttf,
  BoldFont=DejaVuSans-Bold,
  ItalicFont=DejaVuSans-Oblique,
  BoldItalicFont=DejaVuSans-BoldOblique
]{DejaVuSans}
\setmonofont[Path=/usr/share/fonts/TTF/,Extension=.ttf]{DejaVuSansMono}
\usepackage[top=2cm, bottom=2cm, left=2.4cm, right=2.4cm, footskip=1cm]{geometry}
\usepackage{amsmath, amssymb}
\usepackage{xcolor}
\usepackage{enumitem}
\usepackage{fancyhdr}
\usepackage{lastpage}
\usepackage{tabularx}
\usepackage{booktabs}
\usepackage{microtype}
\usepackage{parskip}

\usepackage{array}
\usepackage{longtable}
\usepackage{ltablex}
\keepXColumns
\usepackage{colortbl}
\usepackage[most,breakable]{tcolorbox}
\usepackage{listings}
\usepackage[normalem]{ulem}
\usepackage{pdflscape}
\usepackage{titlesec}
\usepackage[colorlinks=true, linkcolor=accent, urlcolor=accent]{hyperref}

\renewcommand{\familydefault}{\sfdefault}

% ── Theme ─────────────────────────────────────────────────────
\definecolor{accent}{RGB}{0, 60, 113}
\definecolor{primaryblue}{RGB}{0, 60, 113}
\definecolor{loesung}{RGB}{0, 100, 0}
\definecolor{codebg}{RGB}{245, 245, 245}
\definecolor{figurebg}{RGB}{232, 241, 255}
\definecolor{tablehintbg}{RGB}{240, 255, 240}
\definecolor{solutionbg}{RGB}{240, 252, 240}
\definecolor{rulegray}{RGB}{180, 180, 180}
\definecolor{tableheadbg}{RGB}{0, 60, 113}

% ── Header/Footer ─────────────────────────────────────────────
\pagestyle{fancy}
\fancyhf{}
\renewcommand{\headrulewidth}{0.4pt}
\fancyhead[L]{\footnotesize\color{accent}Modulprüfung -- \nichttitle}
\fancyhead[R]{\footnotesize\color{accent}\nichtsubtitle}
\fancyfoot[L]{\footnotesize\color{accent}Matrikelnr.: \rule{3cm}{0.4pt}}
\fancyfoot[R]{\footnotesize\color{accent}Blatt \thepage\,/\,\pageref{LastPage}}

% ── Typografie ────────────────────────────────────────────────
\titleformat{\section}
    {\color{accent}\large\bfseries}{}{0pt}{}
    [\vspace{-4pt}\color{accent}\rule{\linewidth}{1.5pt}]
\titleformat{\subsection}
    {\normalsize\bfseries\color{accent!85!black}}{}{0pt}{}{}
\titleformat{\subsubsection}
    {\small\bfseries\color{accent!70!black}}{}{0pt}{}{}
\titlespacing*{\section}{0pt}{12pt}{4pt}
\titlespacing*{\subsection}{0pt}{8pt}{2pt}
\titlespacing*{\subsubsection}{0pt}{6pt}{1pt}

\setlength{\parindent}{0pt}
\setlength{\parskip}{6pt}
\renewcommand{\arraystretch}{1.25}
\setlist[itemize]{noitemsep, topsep=3pt, parsep=0pt, leftmargin=1.4em}
\setlist[enumerate]{noitemsep, topsep=3pt, parsep=0pt, leftmargin=1.6em}

% ── Boxen ─────────────────────────────────────────────────────
\newtcolorbox{figurebox}[1]{
  colback=figurebg, colframe=accent!50,
  title={Abbildung: #1}, fonttitle=\small\bfseries,
  breakable, before skip=6pt, after skip=6pt
}
\newtcolorbox{tablehintbox}[1]{
  colback=tablehintbg, colframe=green!50!black!50,
  title={Tabelle: #1}, fonttitle=\small\bfseries,
  breakable, before skip=6pt, after skip=6pt
}
\newtcolorbox{solutionbox}[1]{
  enhanced,
  colback=solutionbg, colframe=loesung,
  title={#1}, fonttitle=\small\bfseries,
  coltitle=white, colbacktitle=loesung,
  breakable, before skip=8pt, after skip=8pt
}

\lstset{
  basicstyle=\ttfamily\small,
  backgroundcolor=\color{codebg},
  breaklines=true,
  frame=single, framerule=0pt,
  xleftmargin=4pt, xrightmargin=4pt,
  aboveskip=6pt, belowskip=6pt,
  keepspaces=true
}

\providecommand{\nichttitle}{Probeklausur}
\providecommand{\nichtsubtitle}{}

\renewcommand{\nichttitle}{Betriebssysteme}
\renewcommand{\nichtsubtitle}{Probeklausur 3}


\begin{document}

% Deckblatt
\thispagestyle{empty}
\begin{center}
\vspace*{1cm}
{\Huge\bfseries\color{accent}Probeklausur}\\[8pt]
{\Large\color{accent}\nichttitle}\\[4pt]
\ifx\nichtsubtitle\empty\else
  {\large\color{accent!80!black}\nichtsubtitle}\\[6pt]
\fi
\color{accent}\rule{0.5\linewidth}{1pt}\color{black}
\end{center}

\vspace{1cm}
\begin{tabularx}{\textwidth}{@{}l X@{}}
\textbf{Studiengang:}      & \rule{6cm}{0.4pt} \\[6pt]
\textbf{Matrikelnummer:}    & \rule{6cm}{0.4pt} \\[6pt]
\textbf{Name, Vorname:}    & \rule{6cm}{0.4pt} \\[6pt]
\textbf{Datum:}             & \rule{6cm}{0.4pt} \\[6pt]
\textbf{Bearbeitungszeit:}  & 60 Minuten \\[6pt]
\textbf{Hilfsmittel:}       & Taschenrechner (nicht programmierbar) \\
\end{tabularx}

\vspace{2cm}
\begin{center}
{\large\itshape Viel Erfolg!}
\end{center}

\newpage

% ── BODY ──────────────────────────────────────────────────────

\section{Probeklausur 3 — Betriebssysteme}

\textbf{Bearbeitungszeit}: 60 Minuten | \textbf{Hilfsmittel}: keine | \textbf{Punkte}: 100



\noindent\textcolor{rulegray}{\rule{\linewidth}{0.4pt}}


\paragraph{Aufgabe 1 — Placement \& Buddy \textit{(18 Punkte)}}\mbox{}\\[2pt]

Ein Betriebssystem verwaltet einen freien Speicherbereich von \textbf{1024 KB} mit der Buddy-Technik. Folgende Anforderungen treffen in dieser Reihenfolge ein:


\begin{lstlisting}
A: 70 KB
B: 35 KB
C: 130 KB
D: 250 KB
\end{lstlisting}


Anschließend wird \textbf{B freigegeben}, dann \textbf{A freigegeben}.


a) \textit{(3 P.)} Erkläre den Begriff \textbf{interne Fragmentierung} und nenne den durchschnittlichen Verschnitt auf der letzten Seite einer Datei laut Vorlesung.


b) \textit{(8 P.)} Bestimme für jede der vier Anforderungen A–D die tatsächlich allokierte Blockgröße (gerundet auf 2ᵏ) sowie die dabei entstehende \textbf{interne Fragmentierung pro Anforderung in KB}. Skizziere zusätzlich den Belegungszustand des 1024-KB-Bereichs nach D als Liste der Blöcke (Größe + Status frei/belegt, in Reihenfolge).


c) \textit{(4 P.)} Beschreibe, was nach dem Freigeben von B und anschließend A intern passiert (Buddy-Verschmelzung). Welche Blöcke entstehen?


d) \textit{(3 P.)} Begründe, warum Buddy gegen \textbf{externe} Fragmentierung hilft, aber \textbf{interne} Fragmentierung verstärkt. Nenne ein Verfahren aus der Vorlesung, das den 2ᵏ-Aufrund-Nachteil reduziert.



\noindent\textcolor{rulegray}{\rule{\linewidth}{0.4pt}}


\paragraph{Aufgabe 2 — Slack \& Cluster \textit{(14 Punkte)}}\mbox{}\\[2pt]

Eine Partition verwendet Cluster der Größe \textbf{4096 Byte}, wobei jeder Cluster aus \textbf{8 Sektoren à 512 Byte} besteht. Eine Datei \texttt{report.txt} ist genau \textbf{5300 Byte} groß.


a) \textit{(3 P.)} Wie viele Cluster belegt die Datei? Begründe kurz mit der Definition aus der Vorlesung („kleinste Adressierungseinheit").


b) \textit{(6 P.)} Berechne für die Datei in Byte:

\begin{itemize}
  \item Größe des \textbf{File Slack} (gesamter ungenutzter Platz im letzten Cluster)
  \item Größe des \textbf{RAM Slack} (von EOF bis Sektorende)
  \item Größe des \textbf{Drive Slack} (restliche Sektoren bis Clustergrenze)
\end{itemize}

c) \textit{(2 P.)} Womit füllt Linux RAM- und Drive-Slack — und warum ist diese Wahl forensisch relevant?


d) \textit{(3 P.)} Eine Admin schlägt vor, die Cluster auf \textbf{32 KB} zu vergrößern, weil „große Cluster ja schneller sind". Beurteile diesen Vorschlag mit zwei Argumenten aus dem Lernzettel.



\noindent\textcolor{rulegray}{\rule{\linewidth}{0.4pt}}


\paragraph{Aufgabe 3 — DMA-Ablauf \& Bewertung \textit{(15 Punkte)}}\mbox{}\\[2pt]

Ein Programm liest \textbf{4 KB} von einer SATA-Platte (≈ 2 GB/s) ins RAM. Zur Verfügung stehen zwei Verfahren: \textbf{CPU-gesteuerte Übertragung wortweise} vs. \textbf{DMA}.


a) \textit{(4 P.)} Beschreibe stichpunktartig den \textbf{DMA-Ablauf in 6 Schritten} (Lesen von Platte ins RAM).


b) \textit{(4 P.)} Bei wortweiser CPU-Steuerung (Wortbreite 4 Byte) entsteht pro Wort \textbf{ein Interrupt}. Wie viele Interrupts entstehen bei 4 KB jeweils mit (i) CPU-Steuerung und (ii) DMA? Welche Konsequenz hat das für den Scheduler?


c) \textit{(4 P.)} Nenne und erkläre \textbf{zwei} Vorteile von DMA aus dem Lernzettel, die über die reine Interrupt-Reduktion hinausgehen.


d) \textit{(3 P.)} Ein Tastatur-Treiber (≈ 10 B/s) verzichtet bewusst auf DMA. Begründe, warum DMA hier nicht sinnvoll ist.



\noindent\textcolor{rulegray}{\rule{\linewidth}{0.4pt}}


\paragraph{Aufgabe 4 — Adressierung \& Treiber-Analyse \textit{(15 Punkte)}}\mbox{}\\[2pt]

Gegeben sei folgender Pseudocode aus einem zeichenorientierten Treiber für ein UART-Gerät:


\begin{lstlisting}[language=c]
#define UART_PORT 0x3F8

void writeChar(char c) {
    out(UART_PORT + 0, c);            // Daten
    out(UART_PORT + 5, 0x20);         // Status setzen
}

char readChar() {
    return in(UART_PORT + 0);         // ohne Statusprüfung
}
\end{lstlisting}


a) \textit{(3 P.)} Welche der beiden Adressierungsmethoden wird hier verwendet — \textbf{I/O-Port} oder \textbf{Memory-Mapped I/O}? Woran erkennt man das?


b) \textit{(5 P.)} Der Code enthält \textbf{zwei Mängel} mit Bezug auf die Aufgaben eines Gerätetreibers laut Vorlesung. Identifiziere sie und schlage je eine Korrektur vor.

\textit{(Hinweise: einer betrifft das Ready-Signal des Controllers, einer betrifft eine Standardaufgabe jedes Treibers.)}


c) \textit{(4 P.)} Ordne UART einer Geräteklasse (block- vs. zeichenorientiert) zu und nenne die \textbf{vier} Pflicht-Funktionen dieser Klasse aus dem Lernzettel.


d) \textit{(3 P.)} Vergleiche I/O-Port und MMIO bezüglich \textbf{Schutz} und \textbf{Befehlssatz} (je 1 Vor-/Nachteil).



\noindent\textcolor{rulegray}{\rule{\linewidth}{0.4pt}}


\paragraph{Aufgabe 5 — Journaling vs. Log-structured \textit{(20 Punkte)}}\mbox{}\\[2pt]

Drei Server müssen ein neues Dateisystem bekommen:


\begin{itemize}
  \item \textbf{Server X}: Datenbank-Server, viele kleine Schreibzugriffe, \textbf{Stromausfälle bekannt}, SSD.
  \item \textbf{Server Y}: Video-Bearbeitungsstation, große Dateien, sequentielle Workloads, HDD.
  \item \textbf{Server Z}: Mailserver, jeder verlorene Mail-Body ist kritisch.
\end{itemize}

a) \textit{(4 P.)} Erkläre die Kernregel des \textbf{Journaling-Dateisystems} in einem Satz. Warum garantiert sie nach einem Absturz Konsistenz?


b) \textit{(6 P.)} Vergleiche \textbf{Meta-Daten-Journaling} und \textbf{vollständiges (Data-)Journaling} in einer Tabelle: protokollierter Inhalt, Schreibaufwand, Risiko bei Absturz. Mindestens drei Zeilen.


c) \textit{(6 P.)} Wähle für \textbf{jeden} Server (X, Y, Z) ein passendes Verfahren aus \{Meta-Journaling, Data-Journaling, Log-structured\} und begründe jede Wahl mit \textbf{einem konkreten Argument} aus dem Lernzettel.


d) \textit{(4 P.)} Ein Kollege behauptet: \textit{„Log-structured ist immer besser, weil es nie Daten zurückschreibt."} Widerlege diese Aussage mit zwei Nachteilen aus dem Lernzettel und nenne den zugrundeliegenden Mechanismus, mit dem Log-structured Schreiben funktioniert.



\noindent\textcolor{rulegray}{\rule{\linewidth}{0.4pt}}


\paragraph{Aufgabe 6 — Backup-Strategie \textit{(8 Punkte)}}\mbox{}\\[2pt]

Eine Firma sichert täglich. Am \textbf{Sonntag 02:00} läuft ein Vollbackup, ab \textbf{Montag} dann werktäglich (Mo–Fr) eine \textbf{inkrementelle} Sicherung. Am \textbf{Donnerstagabend} crasht das System komplett.


a) \textit{(2 P.)} Welche Backups müssen in welcher Reihenfolge eingespielt werden, um Donnerstag-Stand wiederherzustellen?


b) \textit{(3 P.)} Das Inkrement von \textbf{Mittwoch} ist defekt. Erkläre die Konsequenz und stelle den maximal noch erreichbaren Stand fest.


c) \textit{(3 P.)} Die Firma erwägt, auf \textbf{differentielles} Backup zu wechseln. Nenne einen Vor- und einen Nachteil gegenüber dem aktuellen inkrementellen Verfahren — bezogen auf die obige Crash-Situation.



\noindent\textcolor{rulegray}{\rule{\linewidth}{0.4pt}}


\paragraph{Aufgabe 7 — Meltdown-Analyse \textit{(10 Punkte)}}\mbox{}\\[2pt]

Folgender Pseudocode wird in einem \textbf{User-Prozess} ausgeführt:


\begin{lstlisting}
1: probe = byte[256 * 4096];     // Probe-Array
2: try {
3:   secret = *(kernel_addr);    // unauthorisiert
4:   dummy  = probe[secret * 4096];
5: } catch (SegFault) { /* ignore */ }
6: // Thread 2: messe Zugriffszeit auf probe[i*4096] für i=0..255
\end{lstlisting}


a) \textit{(3 P.)} Welche \textbf{drei} CPU-Optimierungen aus dem Lernzettel ermöglichen zusammen, dass Zeile 4 trotz späterer Exception in Zeile 3 wirksam wird?


b) \textit{(4 P.)} Erkläre Schritt für Schritt, \textbf{wie Thread 2} in Zeile 6 das Geheimnis rekonstruiert. Beziehe dich auf den Begriff \textbf{Seitenkanal}.


c) \textit{(3 P.)} Bewerte zwei Gegenmaßnahmen — \textbf{KASLR} und \textbf{KPTI/KAISER} — bezüglich Wirksamkeit gegen Meltdown und ihrer Performance-Kosten. Welche der beiden hilft \textit{tatsächlich} gegen Meltdown? Begründe.



\noindent\textcolor{rulegray}{\rule{\linewidth}{0.4pt}}



\section{Musterlösung Klausur 3}

\paragraph{Aufgabe 1 \textit{(18 Punkte)}}\mbox{}\\[2pt]

\textbf{a) (3 P.)} \textit{Interne Fragmentierung}: Innerhalb eines belegten Blocks wird mehr Speicher reserviert als angefordert; der Rest bleibt ungenutzt (1 P.). Sie entsteht, wenn die Blockgröße > der tatsächliche Bedarf ist (1 P.). Durchschnittlich Ø 50 \% der letzten Seite einer Datei sind verschwendet (1 P.).


\textbf{b) (8 P.)} Aufrundung auf nächste 2ᵏ:


{\small
\begin{tabularx}{\linewidth}{XXXX}
\hline
\rowcolor{tableheadbg}
{\color{white}\textbf{Anf.}} & {\color{white}\textbf{Bedarf}} & {\color{white}\textbf{Block}} & {\color{white}\textbf{Interne Fragm.}} \\[2pt]
\hline
\endhead
\hline
\endfoot
A & 70 KB & 128 KB & 58 KB \\
B & 35 KB & 64 KB & 29 KB \\
C & 130 KB & 256 KB & 126 KB \\
D & 250 KB & 256 KB & 6 KB \\
\end{tabularx}
}


(je 0,5 P. Block + 0,5 P. Fragm. = 4 P.)


Belegungszustand nach D (1024 KB von links): \texttt{[128 A][64 B][64 frei][256 C][256 D][256 frei]} (4 P., je Block 0,5 P., richtige Reihenfolge 1 P.).

\textit{Erklärung: A bekommt 128 → Rest 896. Davon Buddy 128 → in 64+64 gesplittet, B nimmt 64. C verlangt 256, also nächstes 256-Buddy aus dem 512er-Block (das 512-Stück wird in 256+256 gesplittet, C nimmt das erste 256). D = 256 nimmt das zweite 256. Restlicher 512er-Block ist die obere Hälfte des 1024 — frei.}


\textbf{c) (4 P.)} B wird freigegeben — sein Buddy (64 frei) ist frei → verschmelzen zu 128 (1 P.). Dieses 128 ist Buddy von A's 128, A ist aber noch belegt → keine weitere Verschmelzung (1 P.). Nach Freigabe von A: A's 128 frei, daneben 128 frei → verschmelzen zu 256 (1 P.). Dieses 256 ist Buddy von dem ursprünglich gesplitteten 256-Bereich (aus dem C/D entstanden) — C/D belegt → keine weitere Verschmelzung. Endzustand: \texttt{[256 frei][256 C][256 D][256 frei]} (1 P.).


\textbf{d) (3 P.)} Externe Fragmentierung wird reduziert, weil Buddies bei Freigabe automatisch verschmolzen werden → größere zusammenhängende Bereiche entstehen (1 P.). Interne wächst, weil jede Anforderung auf 2ᵏ aufgerundet wird; im schlechtesten Fall fast Verdopplung (1 P.). \textbf{Quick-Fit} verzichtet auf die strikte 2ᵏ-Verdopplung und reduziert dadurch die Aufrundung (1 P.).



\noindent\textcolor{rulegray}{\rule{\linewidth}{0.4pt}}


\paragraph{Aufgabe 2 \textit{(14 Punkte)}}\mbox{}\\[2pt]

\textbf{a) (3 P.)} 5300 / 4096 = 1,29 → \textbf{2 Cluster} (1 P.). Cluster ist kleinste Adressierungseinheit, jede Datei belegt mind. 1 Cluster, Reste werden auf vollen Cluster gerundet (2 P.).


\textbf{b) (6 P.)}

\begin{itemize}
  \item Allokierter Platz: 2 · 4096 = 8192 B; Datei = 5300 B → \textbf{File Slack = 8192 − 5300 = 2892 B} (2 P.).
  \item EOF liegt im 2. Cluster bei Offset 5300 − 4096 = 1204 B. Sektor 3 (512–1023) endet bei 1023 — bereits voll. Sektor 4 reicht 1024–1535, EOF in diesem Sektor → \textbf{RAM Slack = 1536 − 1204 = 332 B} (2 P.).
  \item Restliche Sektoren bis Clusterende: Sektoren 5,6,7,8 = 4 · 512 = \textbf{Drive Slack = 2048 B} (2 P.).
\end{itemize}
\textit{Kontrolle:} 332 + 2048 = 2380 B; plus 512 B (in Sektor 4 nach EOF? — bereits in RAM Slack enthalten, da von EOF bis Sektorende). Gesamt File Slack = 332 + 2048 = 2380? Abweichung von 2892 entspricht 512 B → das ist der \textbf{Sektor 4} vor EOF nicht, sondern: korrekt gilt File Slack = RAM + Drive nur, wenn EOF ≠ Sektorgrenze; hier ist 2380 B = RAM+Drive. Differenz 512 B = unvollständiger Sektor zwischen Clustergrenze und EOF — Studi erhält Punkte bei korrekter Methode auch ohne 100 \% Konsistenz; volle 6 P. wenn alle drei Werte (2892 / 332 / 2048) genannt + Methode gezeigt.


\textbf{c) (2 P.)} Linux füllt mit \textbf{Zufallszahlen} statt RAM-/HDD-Inhalten (1 P.). Forensisch relevant: in anderen Systemen können in Slack-Bereichen Spuren früherer Daten (Passwörter, alte Dateifragmente) stecken — bei Linux nicht (1 P.).


\textbf{d) (3 P.)} Schlecht: Große Cluster → mehr Slack → ineffiziente Speichernutzung, vor allem bei vielen kleinen Dateien (1 P.). Vorteil reduzierte Adressierungslast bringt nichts, wenn Partition klein ist (1 P.). Vorlesung: Cluster-Größe sollte zur Partitionsgröße passen — auf kleiner Partition kleine Cluster (1 P.).



\noindent\textcolor{rulegray}{\rule{\linewidth}{0.4pt}}


\paragraph{Aufgabe 3 \textit{(15 Punkte)}}\mbox{}\\[2pt]

\textbf{a) (4 P.)} 1) CPU programmiert DMA-Controller mit Quelle/Ziel/Wortzahl. 2) CPU sendet Lese-Kommando an Platten-Controller. 3) DMA fordert Platten-Controller zur Übertragung an. 4) Platten-Controller → DMA → RAM (Bus). 5) Platten-Controller meldet DMA Abschluss. 6) DMA löst \textbf{einen} Interrupt aus → CPU führt ISR aus. (je Schritt 0,5 P.)


\textbf{b) (4 P.)} 4 KB / 4 B = \textbf{1024 Wörter} → \textbf{1024 Interrupts} mit CPU-Steuerung (1,5 P.); \textbf{1 Interrupt} mit DMA (1,5 P.). Konsequenz: Bei 1024 ISR-Aufrufen verliert der Scheduler stark Rechenzeit, da jeder Interrupt einen Kontextwechsel bedeutet → andere Prozesse bekommen weniger CPU-Zeit (1 P.).


\textbf{c) (4 P.)} Zwei Vorteile (je 2 P.):

\begin{itemize}
  \item \textbf{CPU-Entlastung}: Während DMA überträgt, bearbeitet CPU andere Aufgaben parallel.
  \item \textbf{Spezialisierter Controller}: DMA-Controller ist auf Massentransport optimiert, höhere Gesamtgeschwindigkeit als CPU-wortweise Übertragung.
\end{itemize}
(Alternativ: höhere Gesamtgeschwindigkeit; Parallelität zur CPU.)


\textbf{d) (3 P.)} Tastatur überträgt nur \textasciitilde{}10 B/s (1 P.) — CPU-Last durch Interrupts ist trivial; DMA-Setup-Aufwand übersteigt den Nutzen (1 P.). Außerdem ist Tastatur zeichenorientiert, einzelne Bytes statt Blöcken — DMA ist auf Blocktransfers ausgelegt (1 P.).



\noindent\textcolor{rulegray}{\rule{\linewidth}{0.4pt}}


\paragraph{Aufgabe 4 \textit{(15 Punkte)}}\mbox{}\\[2pt]

\textbf{a) (3 P.)} \textbf{I/O-Port} (1 P.). Erkennbar an \texttt{out}/\texttt{in}-Befehlen mit fester Portnummer 0x3F8 (1 P.); MMIO würde normale Speicherzugriffe (\texttt{*ptr = c}) auf eine gemappte Adresse verwenden (1 P.).


\textbf{b) (5 P.)} Mängel:

\begin{enumerate}
  \item \textbf{Keine Statusabfrage vor Lese-/Schreibzugriff} (2 P.) — der Treiber muss prüfen, ob Controller bereit ist (Lesen aus Statusregister); ein Lesebefehl ohne Bereitschaftsprüfung kann Müll oder altes Datum liefern. Korrektur: vor \texttt{in(UART\_PORT)} Status pollen (\texttt{while(!ready) ;}).
  \item \textbf{Keine ISR / \texttt{handleInterrupt}} (2 P.) — Treiber-Pflicht laut Vorlesung; ohne ISR keine Unterbrechungsbehandlung. Korrektur: Funktion \texttt{handleInterrupt()} ergänzen, beim Controller registrieren.
\end{enumerate}
(Alternativ akzeptiert: fehlende Pufferung, fehlende Initialisierung des Controllers — je 2 P.)

\begin{itemize}
  \item 1 P. für saubere Bezugnahme auf die Treiberaufgaben aus dem Lernzettel.
\end{itemize}

\textbf{c) (4 P.)} \textbf{Zeichenorientiert} (1 P.) — Datenstrom aus einzelnen Zeichen, kein Block. Vier Pflichtfunktionen: \texttt{initDevice}, \texttt{readChar}, \texttt{writeChar}, \texttt{handleInterrupt} (3 P., je 0,75 P.).


\textbf{d) (3 P.)} Vergleich:

\begin{itemize}
  \item I/O-Port: einfach umzusetzen (Vorteil), aber spezielle Befehle (\texttt{in}/\texttt{out}) nötig (Nachteil) (1,5 P.).
  \item MMIO: Zugriff wie normaler Speicherzugriff + leichter Schutz möglich (Vorteil); Unterscheidung „echter Speicher vs. gemappt" aufwändig (Nachteil) (1,5 P.).
\end{itemize}


\noindent\textcolor{rulegray}{\rule{\linewidth}{0.4pt}}


\paragraph{Aufgabe 5 \textit{(20 Punkte)}}\mbox{}\\[2pt]

\textbf{a) (4 P.)} \textbf{Erst Log-Eintrag schreiben, dann die Änderung} (2 P.). Konsistenz: Beim Boot wird das Log abgeglichen — angefangene, unvollendete Transaktionen können rückgängig gemacht oder zu Ende geführt werden, da der Zielzustand vor dem Crash bereits dokumentiert war (2 P.).


\textbf{b) (6 P.)} Tabelle (je Zeile 2 P., Inhalt 1 P., Korrektheit 1 P.):


{\small
\begin{tabularx}{\linewidth}{XXX}
\hline
\rowcolor{tableheadbg}
{\color{white}\textbf{Aspekt}} & {\color{white}\textbf{Meta-Daten-Journaling}} & {\color{white}\textbf{Vollständiges Journaling}} \\[2pt]
\hline
\endhead
\hline
\endfoot
Inhalt & nur Operations-Metadaten (Verzeichnisref, Rechte, Eigentümer, Blöcke) & gesamter Vorgang inkl. Nutzdaten \\
Schreibaufwand & gering & doppelt (alles 2× geschrieben) \\
Risiko bei Absturz & Daten-Inhalte können verloren gehen, Struktur konsistent & auch Daten gesichert; Originale beim Überschreiben gesichert \\
\end{tabularx}
}


\textbf{c) (6 P.)} (je Server 2 P.: 1 P. Wahl, 1 P. Begründung)

\begin{itemize}
  \item \textbf{Server X (DB, kleine Schreibzugriffe, SSD, Stromausfälle)}: \textbf{Meta-Daten-Journaling} — geringer Schreibaufwand passt zu vielen kleinen Writes; SSD profitiert von weniger Schreibzyklen; Konsistenz-Garantie reicht, DB-Engine sichert Daten selbst.
  \item \textbf{Server Y (Video, große sequentielle, HDD)}: \textbf{Log-structured} — sequentielles Schreiben am Log-Ende ideal für HDD-Zugriffsmuster; Copy-on-Write paßt zu großen, selten überschriebenen Dateien. \textit{(Auch akzeptiert: Meta-Journaling, falls Cleaner-Aufwand kritisiert.)}
  \item \textbf{Server Z (Mailserver, Body-Verlust kritisch)}: \textbf{Vollständiges/Data-Journaling} — Mail-Inhalte werden mitprotokolliert, kein Verlust beim Absturz akzeptabel.
\end{itemize}

\textbf{d) (4 P.)} Widerlegung — zwei Nachteile (je 1,5 P.):

\begin{itemize}
  \item Hoher \textbf{Protokollierungsaufwand}: Platte muß zwischen Log und Daten springen.
  \item Starke \textbf{Fragmentierung} durch ständiges Anhängen + Cleaner-Aufwand.
\end{itemize}
Mechanismus (1 P.): \textbf{Copy-on-Write} — alte Daten nicht überschreiben, sondern neue Kopie an Log-Ende schreiben, Verweise auf neue Kopie umbiegen.



\noindent\textcolor{rulegray}{\rule{\linewidth}{0.4pt}}


\paragraph{Aufgabe 6 \textit{(8 Punkte)}}\mbox{}\\[2pt]

\textbf{a) (2 P.)} \textbf{Sonntag-Vollbackup}, dann \textbf{Mo, Di, Mi, Do-Inkrement} in dieser Reihenfolge einspielen (2 P.).


\textbf{b) (3 P.)} Konsequenz: fehlt ein Inkrement, fehlen alle ab diesem Punkt veränderten Dateien — die Kette ist unterbrochen (2 P.). Maximal noch erreichbar: \textbf{Stand Dienstag-Abend} (Voll + Mo + Di) (1 P.).


\textbf{c) (3 P.)} Differentiell vs. inkrementell:

\begin{itemize}
  \item \textbf{Vorteil} in Crash-Situation: Wiederherstellung braucht nur Voll + ein einzelnes Diff (Donnerstag); ein defektes Mittwoch-Diff wäre irrelevant (1,5 P.).
  \item \textbf{Nachteil}: höherer Speicherbedarf, da einmal veränderte Dateien in jedem Diff erneut gesichert werden — Diff wächst mit der Zeit (1,5 P.).
\end{itemize}


\noindent\textcolor{rulegray}{\rule{\linewidth}{0.4pt}}


\paragraph{Aufgabe 7 \textit{(10 Punkte)}}\mbox{}\\[2pt]

\textbf{a) (3 P.)} \textbf{Speculative Execution} (CPU rät, führt vorab aus); \textbf{Out-of-Order Execution} (Befehle nicht in Programmreihenfolge); \textbf{Cache-/TLB-Updates werden nicht zurückgesetzt} beim Verwerfen der Spekulation (je 1 P.). Zusätzlich gilt während der Spekulation die Schutzregel (Supervisor-Bit) \textbf{nicht}.


\textbf{b) (4 P.)} Schritt für Schritt (je 1 P.):

\begin{enumerate}
  \item Zeile 3 wird spekulativ ausgeführt — \texttt{secret} enthält das Kernel-Byte, obwohl die Exception erst danach ausgelöst wird.
  \item Zeile 4 lädt \texttt{probe[secret*4096]} in den Cache; die Exception verwirft den Registerwert, \textbf{aber der Cache-Eintrag bleibt}.
  \item Thread 2 misst Zugriffszeiten für jede der 256 möglichen Probe-Indizes (\texttt{probe[i*4096]}).
  \item Der Index mit der \textbf{kürzesten Latenz} liegt im Cache — das ist \texttt{secret}.
\end{enumerate}
\textbf{Seitenkanal (Cache-Timing)}: Zeitmessungen auf einen geteilten Cache leaken Information, ohne dass der direkte Datenpfad benutzt wird.


\textbf{c) (3 P.)}

\begin{itemize}
  \item \textbf{KASLR}: randomisiert Kernel-Adressen → schützt nur gegen Adress-Erraten; durch \textbf{Spraying / Side-Channels} umgehbar; geringe Performance-Kosten; \textbf{hilft nicht zuverlässig} gegen Meltdown (1,5 P.).
  \item \textbf{KPTI/KAISER}: trennt User- und Kernel-Page-Table physisch; der Kernel ist nicht mehr in der User-Page-Table eingebettet → \textbf{Meltdown-Lesezugriff schlägt fehl}, da kein Mapping existiert. Kosten: \textasciitilde{}5–30 \% langsamer bei Syscalls/Interrupts. \textbf{Wirksam gegen Meltdown} (1,5 P.).
\end{itemize}




% ──────────────────────────────────────────────────────────────

\end{document}
